% \yc{This is the same as the RLC paper from last eyar. Can we rewrite/modify so that it doesnt look like we copied. Also mention how we are different from Melynk, that they have a direct parametrization of reward with policy, where we only have access to samples and use empirical approximation of distribution i.e our approiach is more general where rewards are generated stochastically and we cannot use the reparameterization trick to express them directly as function of policy parameters} 

% Balancing instruction-following and safety in LLMs remains a central challenge for alignment (Henderson et al., 2017; Dinan et al., 2021; Xu et al., 2021; Thoppilan et al., 2022; Bai et al., 2022a,b; Touvron et al., 2023; Dai et al., 2023). While some safety behaviors can align with user intent (e.g., avoiding bias/toxicity), others require refusals for harmful requests (Dinan et al., 2021; Bai et al., 2022a). Early approaches relied on safety critics/filters and dataset curation to reduce unsafe generations (Xu et al., 2021; Thoppilan et al., 2022; Ziegler et al., 2022). In contrast, early RLHF-style instruction following often used a single reward model to jointly encode helpfulness and safety trade-offs (Ouyang et al., 2022; Bai et al., 2022b), which can be sensitive to annotation ambiguity or tuning choices when objectives conflict (Ouyang et al., 2022; Bai et al., 2022b).

% To better manage this trade-off, later work trained separate models for helpfulness and safety: some combine model outputs directly (Glaese et al., 2022; Mu et al., 2024), while others treat safety as an explicit constraint during optimization (Touvron et al., 2023; Ji et al., 2023). Safe RLHF formalizes this constrained perspective by maximizing a reward model subject to a cost constraint, cast in an MDP/constrained-RL framework (Dai et al., 2023; Altman, 2021), and has influenced subsequent safety-constrained RLHF methods (Liu et al., 2024; Huang et al., 2024; Peng et al., 2025). Alternative formulations pursue preference-based balancing objectives (Rame et al., 2023; Zhang et al., 2024; Wachi et al., 2024; Tan et al., 2025). These approaches primarily operationalize safety through expectation-level constraints on harmfulness; high-confidence variants such as HC-RLHF add probabilistic guarantees (via the Seldonian framework) for satisfying the expected-cost constraint using a held-out safety test (Thomas et al., 2019; Chittepu et al., 2025). Our work complements this constrained-RLHF line by replacing expectation-based safety control with distributional dominance constraints on the entire cost distribution, targeting tail-risk control rather than only average-cost control.

% \aj{----------------------------------------------------}

% A growing line of work argues that expectation- or sample-level objectives can miss distributional desiderata, and therefore uses stochastic dominance as an explicit ordering criterion, sometimes operationalized through optimal transport (OT) formulations for tractable learning and matching \citep{melnyk2024distributional,farajzadeh2025imitation,cen2024beyond}. In parallel, optimization research studies how to enforce stochastic-dominance constraints within stochastic programs—via dual/Lagrangian viewpoints and surrogate methods—to make dominance-constrained optimization computationally practical \citep{dentcheva2003optimization,dai2023learning,cen2024beyond}. In safe RLHF, existing approaches decouple helpfulness and safety by learning separate reward and cost models and then optimize reward subject to (high-confidence) cost constraints, typically via Lagrangian-style updates and confidence-bound-based selection/testing \citep{dai2023safe,chittepu2025reinforcement}. Motivated by these threads, our work adopts the constrained safe-RLHF perspective but targets distributional safety control by enforcing that the learned policy’s cost distribution is stochastically smaller than that of the reference policy,
% % first-order stochastic dominance on the learned cost distribution relative to a reference policy,
% leveraging OT-based relaxations/algorithms to make the dominance constraint differentiable and optimizable. Our key difference is replacing scalar (high-confidence) cost constraints in safe RLHF with an OT-optimized FSD constraint on the full cost distribution, yielding distributional safety guarantees not targeted by prior safe-RLHF formulations.

% Balancing instruction-following and safety in LLMs remains a central challenge for alignment \citep{henderson2018ethical,dinan2021anticipating,xu2020recipes,thoppilan2022lamda,bai2022training,bai2022constitutional,touvron2023llama,dai2023safe}.
% % (Henderson et al., 2017; Dinan et al., 2021; Xu et al., 2021; Thoppilan et al., 2022; Bai et al., 2022a,b; Touvron et al., 2023; Dai et al., 2023).
% While some safety behaviors can align with user intent (e.g., avoiding bias/toxicity), others require refusals for harmful requests \citep{dinan2021anticipating,bai2022training}.
% % (Dinan et al., 2021; Bai et al., 2022a).
% Early approaches relied on safety critics/filters and dataset curation to reduce unsafe generations \citep{xu2020recipes,thoppilan2022lamda,ziegler2022adversarial}.
% % (Xu et al., 2021; Thoppilan et al., 2022; Ziegler et al., 2022).
% In contrast, early RLHF-style instruction following often used a single reward model to jointly encode helpfulness and safety trade-offs \citep{ouyang2022training,bai2022constitutional},
% % (Ouyang et al., 2022; Bai et al., 2022b),
% which can be sensitive to annotation ambiguity or tuning choices when objectives conflict \citep{ouyang2022training,bai2022constitutional}.
% % (Ouyang et al., 2022; Bai et al., 2022b).

% To better manage this trade-off, later work trained separate models for helpfulness and safety: some combine model outputs directly \citep{glaese2022improving,mu2024rule},
% % (Glaese et al., 2022; Mu et al., 2024),
% while others treat safety as an explicit constraint during optimization \citep{touvron2023llama,ji2023beavertailsimprovedsafetyalignment}.
% %  (Touvron et al., 2023; Ji et al., 2023).
% Safe RLHF formalizes this constrained perspective by maximizing a reward model subject to a cost constraint, cast in an MDP/constrained-RL framework \citep{dai2023safe,altman2021constrained},
% % (Dai et al., 2023; Altman, 2021),
% and has influenced subsequent safety-constrained RLHF methods \citep{liu2024enhancing,huang2024one,peng2024enhancing}.
% % (Liu et al., 2024; Huang et al., 2024; Peng et al., 2025).
% Alternative formulations pursue preference-based balancing objectives \cite{rame2023rewarded,zhang2024bi,wachi2024stepwise,tan2025equilibrate}.
% % (Rame et al., 2023; Zhang et al., 2024; Wachi et al., 2024; Tan et al., 2025).
% These approaches primarily operationalize safety through expectation-level constraints on harmfulness; high-confidence variants such as HC-RLHF add probabilistic guarantees (via the Seldonian framework) for satisfying the expected-cost constraint using a held-out safety test \citep{thomas2019preventing,chittepu2025reinforcement}.
% % (Thomas et al., 2019; Chittepu et al., 2025).
% Our work complements this constrained-RLHF line by replacing expectation-based safety control with distributional dominance constraints on the entire cost distribution, targeting tail-risk control rather than only average-cost control.

% \aj{----------------------------------------------------}

% A growing line of work argues that expectation- or sample-level objectives can miss distributional desiderata, and therefore uses stochastic dominance as an explicit ordering criterion, sometimes operationalized through optimal transport (OT) formulations for tractable learning and matching \citep{melnyk2024distributional,farajzadeh2025imitation,cen2024beyond}. In parallel, optimization research studies how to enforce stochastic-dominance constraints within stochastic programs—via dual/Lagrangian viewpoints and surrogate methods—to make dominance-constrained optimization computationally practical \citep{dentcheva2003optimization,dai2023learning,cen2024beyond}. In safe RLHF, existing approaches decouple helpfulness and safety by learning separate reward and cost models and then optimize reward subject to (high-confidence) cost constraints, typically via Lagrangian-style updates and confidence-bound-based selection/testing \citep{dai2023safe,chittepu2025reinforcement}. Motivated by these threads, our work adopts the constrained safe-RLHF perspective but targets distributional safety control by enforcing that the learned policy’s cost distribution is stochastically smaller than that of the reference policy,
% % first-order stochastic dominance on the learned cost distribution relative to a reference policy,
% leveraging OT-based relaxations/algorithms to make the dominance constraint differentiable and optimizable. Our key difference is replacing scalar (high-confidence) cost constraints in safe RLHF with an OT-optimized FSD constraint on the full cost distribution, yielding distributional safety guarantees not targeted by prior safe-RLHF formulations.

% Balancing instruction-following and safety in LLMs remains a central challenge for alignment \citep{henderson2018ethical,dinan2021anticipating,xu2020recipes,thoppilan2022lamda,bai2022training,bai2022constitutional,touvron2023llama,dai2023safe}. While some safety behaviors naturally align with user intent (e.g., avoiding bias or toxicity), others require explicit refusals of harmful requests \citep{dinan2021anticipating,bai2022training}. Early approaches mitigated unsafe generations using safety critics, filtering mechanisms, and curated datasets \citep{xu2020recipes,thoppilan2022lamda,ziegler2022adversarial}. In contrast, early RLHF-style instruction-following methods typically employed a single reward model to jointly encode helpfulness and safety trade-offs \citep{ouyang2022training,bai2022constitutional}, a formulation that can be sensitive to annotation ambiguity or tuning choices when these objectives conflict \citep{ouyang2022training,bai2022constitutional}.

% To better manage this trade-off, subsequent work decoupled helpfulness and safety by training separate reward and safety models: some approaches combine model outputs directly \citep{glaese2022improving,mu2024rule}, while others treat safety as an explicit constraint during optimization \citep{touvron2023llama,ji2023beavertailsimprovedsafetyalignment}. Safe RLHF formalizes this constrained perspective by maximizing a reward model subject to a cost constraint within an MDP/constrained-RL framework \citep{dai2023safe,altman2021constrained}, influencing later safety-constrained RLHF methods \citep{liu2024enhancing,huang2024one,peng2024enhancing}. Alternative formulations instead pursue preference-based balancing objectives \citep{rame2023rewarded,zhang2024bi,wachi2024stepwise,tan2025equilibrate}. Across these approaches, safety is primarily operationalized through expectation-level constraints on harmfulness; high-confidence variants such as HC-RLHF extend this framework by providing probabilistic guarantees—via the Seldonian framework—that the expected-cost constraint is satisfied using a held-out safety test \citep{thomas2019preventing,chittepu2025reinforcement}.

\aj{Added two more lines}
Standard RLHF typically optimizes LLM behavior using a single preference-based reward signal, which couples helpfulness and harmlessness and can yield either unsafe completions or unhelpful blanket refusals in adversarial or sensitive prompts \citep{ouyang2022training,bai2022constitutional}. Complementary mitigation strategies aim to reduce harmful outputs using safety critics, filtering mechanisms, and curated datasets \citep{xu2020recipes,thoppilan2022lamda,ziegler2022adversarial}.
To make the objectives more controllable, later approaches separate helpfulness and harmlessness signals, e.g., by combining distinct model scores or by using safety as an explicit optimization constraint \citep{glaese2022improving,mu2024rule,touvron2023llama,ji2023beavertailsimprovedsafetyalignment}.
Safe RLHF makes this trade-off explicit by training separate reward and cost models and then maximizing reward under an expected-cost constraint in a constrained-MDP view \citep{dai2023safe,altman2021constrained}. HC-RLHF further addresses statistical uncertainty by instantiating this constraint within the Seldonian framework and adding a separate safety-screening stage based on upper-confidence bounds computed from held-out data \citep{thomas2019preventing,chittepu2025reinforcement}. %For a more detailed literature review on this, please refer to \cite{chittepu2025reinforcement}.

% Standard RLHF typically optimizes LLM behavior using a single preference-based reward signal that conflates helpfulness and harmlessness, which can lead to either unsafe completions or overly conservative blanket refusals on adversarial or sensitive prompts \citep{ouyang2022training,bai2022constitutional}. Safe RLHF makes this trade-off explicit by learning separate reward and cost models and maximizing reward subject to an expected-cost constraint in a constrained-MDP framework \citep{dai2023safe,altman2021constrained}. HC-RLHF further addresses statistical uncertainty by instantiating this constraint within the Seldonian framework and introducing a separate safety-screening stage based on upper-confidence bounds estimated from held-out data \citep{thomas2019preventing,chittepu2025reinforcement}. For a detailed review of this line of work, see \cite{chittepu2025reinforcement}.

% Alignment for LLMs still hinges on balancing instruction-following with safety \citep{henderson2018ethical,dinan2021anticipating,xu2020recipes,thoppilan2022lamda,bai2022training,bai2022constitutional,touvron2023llama,dai2023safe}. Some safety behaviors reinforce user intent (e.g., reducing bias or toxicity), but other cases require refusing harmful requests \citep{dinan2021anticipating,bai2022training}. Earlier systems reduced unsafe outputs through safety critics, filtering pipelines, and dataset curation \citep{xu2020recipes,thoppilan2022lamda,ziegler2022adversarial}. By contrast, early RLHF-style instruction-following approaches commonly used one reward model to represent both helpfulness and safety trade-offs \citep{ouyang2022training,bai2022constitutional}, which can become brittle when annotation ambiguity or tuning decisions sharpen conflicts between the two objectives \citep{ouyang2022training,bai2022constitutional}.

% To address this trade-off, later methods separate helpfulness and safety by learning distinct reward and safety models: some combine model outputs directly \citep{glaese2022improving,mu2024rule}, whereas others impose safety as an explicit optimization constraint \citep{touvron2023llama,ji2023beavertailsimprovedsafetyalignment}. Safe RLHF crystallizes this constrained view by optimizing a reward model under a cost constraint in an MDP/constrained-RL formulation \citep{dai2023safe,altman2021constrained}, and this setup has shaped subsequent safety-constrained RLHF approaches \citep{liu2024enhancing,huang2024one,peng2024enhancing}. Other lines of work focus on preference-based balancing objectives \citep{rame2023rewarded,zhang2024bi,wachi2024stepwise,tan2025equilibrate}. Overall, these methods mostly define safety through expectation-level harmfulness constraints; high-confidence variants such as HC-RLHF augment this with probabilistic guarantees through the Seldonian framework, that a held-out safety test verifies satisfaction of the expected-cost constraint \citep{thomas2019preventing,chittepu2025reinforcement}.

In parallel, a growing line of work argues that expectation objectives can overlook distributional desiderata, advocating stochastic dominance as an explicit ordering criterion and, in some cases, leveraging optimal transport (OT) formulations for tractable learning and matching \citep{melnyk2024distributional,farajzadeh2025imitation,cen2024beyond}. The closest complementary advances in stochastic optimization study how to enforce stochastic-dominance constraints via dual/Lagrangian characterizations and surrogate methods to make dominance-constrained optimization computationally practical \citep{dentcheva2003optimization,dai2023learning,cen2024beyond}. Building on both the constrained-RLHF and dominance-based optimization perspectives, our work adopts the safe-RLHF framework but replaces expectation-based safety control with a first-order stochastic dominance constraint on the entire cost distribution relative to a reference policy, leveraging OT-based relaxations to render the dominance constraint differentiable and optimizable. Unlike prior Safe RLHF methods that impose scalar (possibly high-confidence) expected-cost constraints \citep{dai2023safe,chittepu2025reinforcement}, our approach enforces an OT-optimized FSD constraint on the full cost distribution, thereby targeting distributional safety guarantees beyond average-cost control.
